% Compiles with revtex4-2 (available by default on Overleaf, template
% "REVTeX"; or `tlmgr install revtex4-2` in a local TeX Live install).
% Target: arXiv physics.comp-ph, cross-list gr-qc.
\documentclass[aps,onecolumn,10pt,nofootinbib]{revtex4-2}

\usepackage{amsmath}
\usepackage{amssymb}
\usepackage{hyperref}

\begin{document}

\title{A tidal-scale vacuum diagnostic for spacetimes with vanishing curvature invariants, and the numerical Einstein tensor as arbiter for exact-solution ambiguities}

\author{Marco Vinicius Amorin de Lima}
\affiliation{Independent researcher, Brazil}
\email{abbeats@gmail.com}

\date{\today}

\begin{abstract}
Numerical relativity and geodesic-integration codes routinely need a local
``is this vacuum?'' diagnostic: compute an effective stress-energy tensor
from the numerical Einstein tensor and compare it against a tolerance scaled
by the local curvature. A common choice of scale is a polynomial curvature
invariant such as the Kretschmann scalar $K$. We show that this choice fails
completely, and silently, in spacetimes with vanishing scalar invariants
(VSI) — the exact plane gravitational wave being the physically relevant
example — because every polynomial curvature invariant is identically zero
there even though the Riemann tensor is not. An invariant-normalized
tolerance collapses to zero and any finite-difference noise floor is then
classified as matter. We give a fix: scale the tolerance by the Riemann
tensor projected onto a local orthonormal tetrad (the physical tidal scale
an extended observer would measure), which remains finite and physically
meaningful in VSI spacetimes while agreeing with the invariant-based scale
in the non-degenerate cases we checked. We demonstrate the failure and the
fix numerically on the exact Brinkmann plane wave. Separately, we report a
general method for resolving genuine algebraic ambiguities encountered when
reconstructing a published exact solution from first principles: when a
sign or factor choice is left undetermined by all known limiting cases, and
the true vacuum/matter content of the solution is independently known, the
numerical field-equation pipeline itself can arbitrate between candidates.
We demonstrate this on the Khan--Penrose colliding-plane-wave metric, where
one of two algebraically valid conformal-factor candidates is rejected
because it requires a matter density seven to nine orders of magnitude
above the finite-difference noise floor in the interaction region. Both
findings emerged from building an interactive general-relativity teaching
tool and are offered here as engineering notes for other codes with
similar diagnostics.
\end{abstract}

\maketitle

\section{Introduction}

A geodesic or curvature code that lets a user probe an arbitrary point of
an arbitrary spacetime needs, at minimum, a way to answer ``is this vacuum,
locally, right here?'' The natural implementation computes the numerical
Einstein tensor $G_{\mu\nu}$ from finite-difference Christoffel symbols (or
from analytic Christoffels differentiated numerically), forms an effective
stress-energy tensor
\begin{equation}
  T^{\text{eff}}_{\mu\nu} = \frac{c^4}{8\pi G}\,G_{\mu\nu},
\end{equation}
and compares its components against a tolerance. Because the raw magnitude
of $T^{\text{eff}}_{\mu\nu}$ depends on the coordinate scale and the
distance from any interesting feature, the tolerance must itself be scaled
by something locally characteristic of the curvature. Two engineering
findings from building such a diagnostic, in the course of constructing an
interactive general-relativity laboratory (CosmoLab), are reported here.

First (Sec.~\ref{sec:vsi}): scaling the tolerance by a polynomial curvature
invariant such as the Kretschmann scalar $K = R_{\alpha\beta\gamma\delta}R^{\alpha\beta\gamma\delta}$
is the obvious choice, and it is what we started with. It fails completely
in spacetimes with vanishing scalar invariants (VSI)
\cite{Pravda2002,Brinkmann1925}, where every such invariant is identically
zero although the Riemann tensor is not. The fix is to scale the tolerance
by the Riemann tensor projected onto a local orthonormal tetrad — the
physical tidal scale, rather than an algebraic contraction that happens to
vanish identically for null-aligned curvature.

Second (Sec.~\ref{sec:arbiter}): when an exact solution is reconstructed
from first principles rather than transcribed from a source text, a genuine
algebraic ambiguity can survive every known limiting case. We report a
general protocol for such situations — let the numerical field equations
arbitrate, using independently known physical content (here, that the
interaction region of a colliding-wave spacetime is vacuum) as the
discriminating criterion — and demonstrate it on the Khan--Penrose metric
\cite{KhanPenrose1971,Griffiths1991}.

Both sections give the failure mode, the fix, and a numerical
demonstration with the actual figures produced by the implementation, so
that the argument can be checked independently of the source code.

\section{Vanishing scalar invariants and the collapse of invariant-normalized vacuum diagnostics}
\label{sec:vsi}

\subsection{Background}

A spacetime has vanishing scalar invariants (VSI) if every polynomial
scalar invariant constructed from the Riemann tensor and its covariant
derivatives is identically zero, even though the Riemann tensor itself is
not \cite{Pravda2002}. The canonical examples are the plane-fronted
gravitational waves with parallel rays (pp-waves), first written in the
form used here by Brinkmann in 1925 \cite{Brinkmann1925}. The underlying
reason is algebraic: the curvature of such a spacetime is of Petrov type N
(or more degenerate), aligned with a null direction, and every scalar
contraction of a null-aligned tensor with itself vanishes for the same
reason a null vector has zero norm without being the zero vector.

The physically relevant instance is the exact plane gravitational wave in
Brinkmann coordinates,
\begin{equation}
  ds^2 = -(dx^0)^2 + dx^2 + dy^2 + dz^2 + H(u,x,y)\,du^2,
  \qquad u = \frac{x^0 - z}{\sqrt{2}},
\end{equation}
\begin{equation}
  H(u,x,y) = A_+(u)\,(x^2 - y^2) + 2A_\times(u)\,xy.
\end{equation}
Because $x^2 - y^2$ and $xy$ are harmonic functions of the transverse
coordinates, $\nabla_\perp^2 H = 0$ and this metric is an \emph{exact}
vacuum solution of the Einstein field equations for \emph{any} profile
$A_+(u), A_\times(u)$ — not a linearized approximation. One direct
consequence, verified numerically below, is that the Ricci scalar $R$ and
the Kretschmann scalar $K = R_{\alpha\beta\gamma\delta}R^{\alpha\beta\gamma\delta}$
both vanish identically, while the Riemann tensor itself does not: no
scalar-invariant detector registers the wave. This is also the underlying
reason a gravitational-wave interferometer measures a distance (geodesic
deviation between test masses) rather than a curvature scalar.

\subsection{The failure mode}
\label{sec:vsi-failure}

A vacuum diagnostic that asks whether $|T^{\text{eff}}_{\mu\nu}|$ is
smaller than a tolerance $\varepsilon$, with $\varepsilon$ scaled by
$\sqrt{K}$, degenerates in a VSI spacetime because $K \equiv 0$: the
tolerance collapses to zero (or to a floor set only by machine epsilon),
and \emph{any} nonzero finite-difference truncation or round-off noise in
the numerically computed $T^{\text{eff}}_{\mu\nu}$ exceeds it. The
diagnostic then reports ``matter'' for an exact vacuum wave. The failure is
silent: no exception, no warning, just a wrong classification, because the
code path that would normally distinguish ``genuinely curved but vacuum''
from ``matter present'' has nothing left to divide by that is not already
zero for algebraic reasons unrelated to whether matter is actually present.

\subsection{The fix: tidal scale from a tetrad-projected Riemann tensor}
\label{sec:vsi-fix}

The fix used here replaces the invariant-based scale with the physical
tidal scale: build a local orthonormal tetrad $\{\hat e_0, \hat e_1, \hat
e_2, \hat e_3\}$ by Gram--Schmidt orthonormalization of a chosen observer
$\hat e_0 \equiv \hat u$ (static, zero-angular-momentum, or Eulerian,
depending on which is available at that point in the chart) against the
coordinate basis, project the Riemann tensor onto that tetrad, and take
\begin{equation}
  \Xi \equiv \max_{\hat a \hat b \hat c \hat d}
    \left| R_{\hat a \hat b \hat c \hat d} \right|,
\end{equation}
the largest tetrad-frame component of the (fully lowered) Riemann tensor —
physically, the maximum tidal stretching an extended observer at that
tetrad would measure. The tolerance is then
\begin{equation}
  \varepsilon = \max\!\left(10^{-4}\,\frac{c^4}{8\pi G}\,\Xi,\ \varepsilon_{\text{floor}}\right),
\end{equation}
with $\varepsilon_{\text{floor}}$ a small absolute floor to avoid division
pathologies at literally flat points. Because $\Xi$ is built from a
tetrad-frame projection of the Riemann tensor rather than an algebraic
self-contraction, it does not depend on the null-alignment cancellation
that drives $K$ to zero: $\Xi$ is finite and nonzero for the plane wave
precisely where $K$ is not. In the one non-VSI case we checked
(Schwarzschild), $\Xi$ and $\sqrt{K}$ agree up to an $O(1)$ factor, so the
fix does not change the diagnostic's behavior where the invariant-based
scale already worked; we do not claim this as a general algebraic identity
for all non-VSI spacetimes, only that the tidal scale remains well-defined
and physically motivated in both regimes, which is what the invariant
scale fails to be in the VSI case.

\subsection{Numerical demonstration}
\label{sec:vsi-demo}

We evaluate a monochromatic wave, $A_+(u) = A_0 \cos(k_u u)$, $A_\times =
0$, with $A_0 = \tfrac{1}{2}k_u^2 h_0$ for a target strain $h_0 = 0.05$ and
wavelength $\lambda_u = 6\times 10^6\,\mathrm{m}$, at a point inside the
wave, $x^0 = 0.2\lambda_u$, $x = 0.8L$, $y = 0.5L$, $z=0$, with
$L = 10^6\,\mathrm{m}$ the finite-difference step scale. The measured
values are
\begin{align}
  R &\approx 1.7\times 10^{-26} \quad (\text{machine-precision noise, i.e. } R = 0 \text{ exactly}), \\
  K &\approx -9.0\times 10^{-44}, \quad \sqrt{|K|} \approx 3.0\times 10^{-22}, \\
  \max|\text{Riemann tensor component}| &\approx 8.6\times 10^{-15} \quad (\text{comparable to } A_0 \approx 2.7\times 10^{-14}\,\mathrm{m^{-2}}), \\
  T^{\text{eff}}(\hat u, \hat u) &\approx 1.9\times 10^{17}\ \mathrm{J/m^3} \quad (\text{finite-difference noise floor}).
\end{align}
An invariant-normalized tolerance would be
$10^{-4}\sqrt{|K|} \approx 3\times 10^{-26}$: eleven orders of magnitude
\emph{below} the measured noise floor of
$T^{\text{eff}}(\hat u,\hat u)$, so the diagnostic would report
``matter present'' for an exact vacuum wave. The tetrad-projected tidal
scale $\Xi$, of the same order as the Riemann-tensor components above
($\sim 10^{-15}$), gives
$\varepsilon = 10^{-4}(c^4/8\pi G)\,\Xi \sim 10^{22}\ \mathrm{J/m^3}$,
comfortably above the noise floor, and the fixed diagnostic correctly
reports vacuum. The same wave, evaluated with the fixed diagnostic,
passes at all sampled points, while the Riemann tensor itself is
nonzero to more than an order of magnitude above the noise floor at the
same point — the curvature is real and the wave is not flat spacetime;
only the polynomial scalar invariants built from it happen to vanish.

\subsection{Recommendation}

Any code with a vacuum/matter diagnostic normalized by a polynomial
curvature invariant has this failure mode latent whenever it is pointed
at a VSI spacetime — plane waves being the case of practical interest,
since they are the standard exact model for a gravitational-wave
spacetime. We recommend scaling such diagnostics by a tetrad-projected
curvature quantity (the tidal scale above, or an equivalent local
observer-dependent norm of the Riemann tensor) rather than by a
polynomial invariant, since the former does not depend on an algebraic
cancellation that has nothing to do with whether matter is present.

\section{The numerical Einstein tensor as arbiter: reconstructing Khan--Penrose from first principles}
\label{sec:arbiter}

\subsection{Background}

The collision of two exact plane gravitational waves is a classic
demonstration of the nonlinearity of general relativity: each wave alone
propagates through flat spacetime harmlessly, but the region where two
colliding waves overlap develops genuine spacetime curvature that a linear
superposition would not produce, and in the Khan--Penrose solution
\cite{KhanPenrose1971} this interaction region collapses to a real
curvature singularity in finite proper time, from pure vacuum gravity with
no matter source. We follow Griffiths' classification of colliding
plane-wave spacetimes \cite{Griffiths1991} throughout.

In collinear-polarization coordinates, with $u,v$ the (Heaviside-clipped,
dimensionless) null phases of the two incoming waves, define
$p \equiv \sqrt{1-u^2}$, $q \equiv \sqrt{1-v^2}$, $W \equiv 1-u^2-v^2$,
$t \equiv uq+vp$. The interaction-region metric takes the form
\begin{equation}
  ds^2 = -F\,(dx^0)^2 + F\,dz^2 + G\,dx^2 + H\,dy^2,
\end{equation}
with
\begin{equation}
  G = \frac{W(1-t)}{1+t}, \qquad
  H = \frac{W(1+t)}{1-t}, \qquad
  GH = W^2.
\end{equation}
The single-wave limiting cases (regions where only one of the two waves
has arrived) fix $G$, $H$, and the leading behavior of $W$ uniquely. They
do \emph{not} fix the conformal factor $F$ uniquely: both
\begin{equation}
  F_\pm = \frac{W^{3/2}}{p\,q\,(pq \pm uv)^2}
\end{equation}
reduce to the correct single-wave limits in the non-interacting regions.
This is a genuine algebraic ambiguity encountered when the metric is
reconstructed from the structure of the field equations rather than
transcribed directly from a source text.

\subsection{Method: the numerical Einstein tensor as arbiter}

The published Khan--Penrose solution is vacuum throughout, including the
interaction region. This is an independently known physical fact about the
solution, not a free parameter — so instead of resolving the sign in
$F_\pm$ by further analytic work, we implemented \emph{both} metrics,
$F_+$ and $F_-$, in the same numerical pipeline (finite-difference
Christoffel symbols, cross-checked against the analytic Christoffel
symbols of the piecewise-diagonal metric to $6\times 10^{-11}$ relative
accuracy in the smooth interior, well inside the interaction region), and
computed the effective stress-energy tensor of each candidate at several
points in the interaction region.

\subsection{Result}

At the sampled points $(u,v) = (0.25,0.25),\ (0.5,0.2),\ (0.6,0.6)$ (with
the focal length scale $a=3\times10^8\,\mathrm{m}$), the two candidates
give:

\begin{center}
\begin{tabular}{c c c c}
\hline\hline
$(u,v)$ & $T^{\text{eff}}(\hat u,\hat u)$ for $F_+$ [J/m$^3$] & $T^{\text{eff}}(\hat u,\hat u)$ for $F_-$ [J/m$^3$] & ratio \\
\hline
$(0.25,0.25)$ & $1.8\times 10^{16}$ & $-3.1\times 10^{25}$ & $1.7\times 10^{9}$ \\
$(0.5,0.2)$   & $3.6\times 10^{16}$ & $-6.8\times 10^{25}$ & $1.9\times 10^{9}$ \\
$(0.6,0.6)$   & $4.0\times 10^{19}$ & $-6.7\times 10^{26}$ & $1.7\times 10^{7}$ \\
\hline\hline
\end{tabular}
\end{center}

The $F_+$ candidate is consistent with vacuum at the same finite-difference
noise-floor level found throughout the rest of the code (compare
Sec.~\ref{sec:vsi-demo}); the $F_-$ candidate requires a matter density
seven to nine orders of magnitude above that floor at every sampled point.
$F_+$, i.e. the $(pq+uv)^2$ choice, is therefore the physically correct
factor, consistent with the published 1971 solution. As an independent
consistency check, the anisotropic (Kasner-type) collapse of the
interaction region, $\sqrt{G}\to 0$ (compression) while $\sqrt{H}\to\infty$
(stretching) as $W\to 0$, satisfies the exact identity
$\sqrt{G}\cdot\sqrt{H} = W$ to $10^{-10}$ relative accuracy for the
selected ($F_+$) candidate, matching the area-factor collapse expected of
the solution independently of which conformal-factor candidate is used
for $F$.

\subsection{Discussion: a general protocol}

The method generalizes beyond this one metric: when an exact solution is
reconstructed from the structure of the field equations rather than copied
from a primary source, and a genuine algebraic ambiguity survives every
known limiting case, the numerical field-equation pipeline itself can
arbitrate between candidates, provided the physical vacuum/matter content
of the true solution is independently known from the literature (here,
that the Khan--Penrose interaction region is vacuum). The discriminating
power of the method rests entirely on that independent knowledge: it
cannot resolve an ambiguity in a spacetime whose true matter content is
itself unknown. Within that scope, the protocol turns a numerical
relativity code's own validation machinery into a referee for its
implementation choices, and we keep both candidates in the code's
regression test suite as a record of the arbitration.

\section{Discussion and recommendations}

Both findings reported here emerged as byproducts of building an
interactive teaching tool for general relativity, CosmoLab, not from a
dedicated numerical-methods research program — but the underlying failure
modes are general to any code that (a) normalizes a vacuum/matter
diagnostic by a polynomial curvature invariant, or (b) reconstructs a
published exact solution from first principles rather than transcription.
We suggest two concrete practices for other geodesic and curvature codes:

\begin{enumerate}
  \item Default vacuum/matter diagnostics to a tetrad-projected tidal
  scale rather than a polynomial curvature invariant, since the latter
  is not a robust normalization whenever the spacetime under test may be
  of Petrov type N or more degenerate (VSI), which is exactly the case
  of a gravitational-wave spacetime — arguably the single most relevant
  case for a relativity code to handle correctly.
  \item When reconstructing an exact solution and a residual algebraic
  ambiguity survives all known limits, use independently known physical
  content plus the numerical field equations as an arbiter, and record
  the rejected candidate as a regression test documenting the decision,
  rather than silently discarding it.
\end{enumerate}

\section{Conclusion}

We reported two engineering findings from the numerical implementation of
exact spacetimes in an interactive relativity laboratory. First, vacuum
diagnostics normalized by polynomial curvature invariants collapse to a
vacuous tolerance in VSI spacetimes such as the exact plane gravitational
wave, misclassifying a genuine vacuum solution as matter; scaling instead
by the tetrad-projected Riemann tensor (the physical tidal scale) fixes
this while leaving non-degenerate cases unaffected. Second, a genuine
algebraic ambiguity encountered while reconstructing the Khan--Penrose
colliding-wave metric from first principles was resolved by using the
numerical Einstein tensor itself as an arbiter against the independently
known vacuum content of the solution, rejecting a candidate that would
require a matter density seven to nine orders of magnitude above the
numerical noise floor. Both results are implemented and covered by
automated regression tests in the CosmoLab code base
\cite{CosmoLab}.

\begin{acknowledgments}
This work is part of the CosmoLab project, an open-source interactive
general-relativity laboratory (DOI: 10.5281/zenodo.21436967). The
scientific claims of that project are subject to an open, adversarial
validation protocol designed to be run against an independent physicist
or reasoning system, available in the project repository.
\end{acknowledgments}

\begin{thebibliography}{99}

\bibitem{Pravda2002}
V.~Pravda, A.~Pravdov\'a, A.~Coley, and R.~Milson,
``All spacetimes with vanishing curvature invariants,''
Class. Quantum Grav. \textbf{19}, 6213 (2002).

\bibitem{Brinkmann1925}
H.~W.~Brinkmann,
``Einstein spaces which are mapped conformally on each other,''
Math. Ann. \textbf{94}, 119 (1925).

\bibitem{KhanPenrose1971}
K.~Khan and R.~Penrose,
``Scattering of two impulsive gravitational plane waves,''
Nature \textbf{229}, 185 (1971).

\bibitem{Griffiths1991}
J.~B.~Griffiths,
\emph{Colliding Plane Waves in General Relativity}
(Oxford University Press, Oxford, 1991).

\bibitem{CosmoLab}
M.~V.~Amorin de Lima,
\emph{CosmoLab: an interactive general-relativity laboratory in the browser},
Zenodo, v1.0.0 (2026), \url{https://doi.org/10.5281/zenodo.21436967}.

\end{thebibliography}

\end{document}
